\section{Compito 1}

\subsection{Descrizione del problema}
Il problema affrontato in questo elaborato consiste nell'ottimizzazione dell'ingaggio di giocatori per una squadra sportiva a budget limitato. Tale formulazione modella il tipico scenario decisionale affrontato dalla dirigenza (in particolare dal \textit{General Manager}) di una franchigia di baseball durante la sessione di mercato dei giocatori svincolati (\textit{free agency}): la necessità di comporre la rosa titolare più competitiva possibile rispettando vincoli finanziari di bilancio e di copertura dei ruoli sul diamante.

In letteratura accademica, il problema è tratto dal testo di riferimento \textit{Introduction to Algorithms} (T.~H.~Cormen, C.~E.~Leiserson, R.~L.~Rivest, C.~Stein), in particolare dal Problema 15-12 (``Signing free-agent baseball players'') presente nel capitolo dedicato alla programmazione dinamica. Nella teoria dell'ottimizzazione combinatoria e della ricerca operativa, tale problema costituisce un'istanza canonica del Multiple-Choice Knapsack Problem (MCKP), variante specializzata del classico problema dello zaino (0/1 Knapsack) in cui gli oggetti sono partizionati in classi disgiunte e mutualmente esclusive.

\subsection{Formulazione Matematica dell'Input}
L'istanza del problema riceve in ingresso un insieme strutturato di dati discreti e deterministici:

\begin{itemize}
    \item \textbf{Budget Totale ($B$):} un intero non negativo $B \in \mathbb{N}$ che rappresenta la disponibilità finanziaria massima spendibile per l'ingaggio complessivo dei componenti della squadra.
    \item \textbf{Numero di Ruoli ($N$):} un intero positivo $N \in \mathbb{N}$ che definisce il numero complessivo di ruoli/reparti tecnici necessari alla formazione della squadra. Nel contesto del baseball si considerano i canonici $N = 9$ ruoli difensivi (es. lanciatore titolare -- \textit{Starting Pitcher}, ricevitore -- \textit{Catcher}, prima base -- \textit{First Base}, ecc.). I ruoli sono indicizzati da $i = 1, 2, \dots, N$.
    \item \textbf{Candidati per Ruolo ($P_i$):} per ciascun ruolo $i \in \{1, \dots, N\}$, è disponibile sul mercato un insieme finito di $P_i$ candidati svincolati, indicizzati con $k \in \{1, 2, \dots, P_i\}$.
    \item \textbf{Parametri dei Giocatori:} ogni candidato $k$ appartenente al ruolo $i$ è univocamente caratterizzato da una coppia di parametri:
    \begin{itemize}
        \item $c_{i,k} \in \mathbb{N}$: costo d'acquisto/ingaggio del giocatore sul mercato dei \textit{free agent}.
        \item $v_{i,k} \in \mathbb{R}^+$: metrica di rendimento tecnico del giocatore, quantificata tramite la statistica sabermetrica \textbf{WAR} (\textit{Wins Above Replacement}).
    \end{itemize}
\end{itemize}

\subsubsection{La Metrica Sabermetrica WAR (Wins Above Replacement)}
Nell'analisi statistica avanzata del baseball (\textit{sabermetrics}), il \textbf{WAR} (\textit{Wins Above Replacement}) rappresenta la metrica sintetica di riferimento per valutare il valore complessivo fornito da un atleta alla propria squadra. 

Il concetto alla base del WAR è il confronto con il cosiddetto \textit{giocatore di livello rimpiazzo} (\textit{replacement-level player}), ovvero un atleta liberamente acquisibile sul mercato al salario minimo contrattuale o promosso dalle leghe minori (Minor League Baseball - MiLB), il cui rendimento incrementale è convenzionalmente posto pari a $0.0$. Un giocatore con un valore $\text{WAR} > 0$ esprime quante vittorie aggiuntive nel corso della stagione (composta da 162 incontri in MLB) la squadra ottiene schierando tale atleta rispetto a una riserva di livello rimpiazzo.

Sebbene nel testo originale di Cormen et al. venga fatto riferimento alla metrica storica \textbf{VORP} (\textit{Value Over Replacement Player}), ideata originariamente per misurare i punti generati in battuta sopra il livello di rimpiazzo, nella moderna dirigenza sportiva e nel nostro modello si adotta il WAR, sua diretta e completa evoluzione. Il WAR integra organicamente ogni aspetto del gioco (battuta, corsa sulle basi, difesa sul campo e lanci).

Sul piano dell'ottimizzazione combinatoria, l'adozione del WAR offre due proprietà teoriche ideali per il paradigma dinamico:
\begin{enumerate}
    \item \textbf{Linearità e Additività:} le vittorie sopra il rimpiazzo sono quantitativamente sommabili tra loro. Il WAR totale della rosa è esattamente dato dalla somma lineare dei contributi dei singoli atleti schierati.
    \item \textbf{Coerenza del Livello Base (Costo Zero):} non ingaggiare alcuno svincolato in un dato ruolo corrisponde a impiegare la riserva già contrattualizzata a costo $0$ con apporto incrementale $\text{WAR} = 0.0$, garantendo la perfetta rispondenza con il caso base dell'algoritmo.
\end{enumerate}

\subsection{Vincoli del Problema}
Una selezione di giocatori è considerata ammissibile se e solo se soddisfa simultaneamente i seguenti due vincoli operativi:

\begin{itemize}
    \item \textbf{Vincolo di Rosa (Esclusività di Ruolo):} per ogni singolo ruolo $i \in \{1, \dots, N\}$, è ammesso l'ingaggio di al massimo un giocatore. È consentito non ingaggiare alcun candidato in un determinato ruolo (mantenendo la riserva d'ufficio presente in rosa a costo $0$ e apporto incrementale nullo).

    Formalmente, introducendo le variabili binarie di decisione:
    \[
    x_{i,k} \in \{0, 1\} \quad \forall i \in \{1, \dots, N\}, \quad \forall k \in \{1, \dots, P_i\}
    \]
    il vincolo si esprime come:
    \[
    \sum_{k=1}^{P_i} x_{i,k} \le 1 \quad \forall i \in \{1, \dots, N\}
    \]

    \item \textbf{Vincolo di Bilancio (Budget Globale):} la somma cumulativa dei costi associati a tutti i giocatori ingaggiati non deve eccedere la soglia di budget $B$ fissata a monte:
    \[
    \sum_{i=1}^N \sum_{k=1}^{P_i} c_{i,k} \cdot x_{i,k} \le B
    \]
\end{itemize}

\subsection{Funzione Obiettivo e Output Richiesto}
Il problema si qualifica come un problema di massimizzazione vincolata.

\begin{itemize}
    \item \textbf{Funzione Obiettivo:} determinare l'assegnamento ottimale delle variabili decisionali $x_{i,k}$ tale da massimizzare il valore cumulativo (competitività totale della rosa) apportato alla squadra:
    \[
    \max \sum_{i=1}^N \sum_{k=1}^{P_i} v_{i,k} \cdot x_{i,k}
    \]
    sotto il rispetto dei vincoli definiti precedentemente.

    \item \textbf{Output Richiesto:} conformemente alle specifiche di progetto, la risoluzione deve produrre:
    \begin{enumerate}
        \item \textbf{Il valore scalare ottimo}, ovvero il massimo WAR complessivo conseguibile senza violare il vincolo finanziario.
        \item \textbf{La configurazione ottima} (la rosa definitiva), descritta mediante l'elenco esatto dei giocatori scelti ruolo per ruolo che realizzano tale rendimento.
    \end{enumerate}
\end{itemize}

\subsection{Dimostrazione dell'Applicabilità della Programmazione Dinamica}
L'applicabilità del paradigma della programmazione dinamica a un problema di ottimizzazione poggia su due presupposti teorici fondamentali:
\begin{enumerate}
    \item la presenza di una \textbf{sottostruttura ottima} all'interno di una soluzione ottima;
    \item l'esistenza di una \textbf{sovrapposizione dei sottoproblemi} (\textit{overlapping subproblems}).
\end{enumerate}

\subsubsection{Definizione dei Sottoproblemi}
Siano $N$ le posizioni da ricoprire (numerate da $1$ a $N$), $P$ il numero di giocatori per ruolo e $X$ il budget nominale disponibile. Per specifica, i costi di ciascun giocatore sono multipli di $100\,000$: introducendo la costante di scala $\text{COST\_UNIT} = 100\,000$, definiamo il budget discreto:
\[
B = \left\lfloor \frac{X}{\text{COST\_UNIT}} \right\rfloor
\]
Analogamente, il costo intero scalato del $j$-esimo giocatore per il ruolo $i$ ($0 \le j < P$) è dato da $c_{i, j} = \left\lfloor \text{players}[i-1, j].cost / \text{COST\_UNIT} \right\rfloor$.

Definiamo formalmente il generico \textbf{sottoproblema} associato alla coppia di parametri $(i, b)$ come:
\begin{quote}
\textit{«Massimizzare il WAR complessivo considerando esclusivamente un sottoinsieme di ruoli tra i primi $i$ (con $0 \le i \le N$), disponendo di una quantità di budget residuo pari al più a $b$ (con $0 \le b \le B$).»}
\end{quote}
Indichiamo con $m[i, b]$ il valore della soluzione ottima di tale sottoproblema.

\subsubsection{Dimostrazione della Sottostruttura Ottima}
Un problema esibisce una sottostruttura ottima se una soluzione ottima al problema contiene al suo interno soluzioni ottime per i relativi sottoproblemi. Dimostriamo tale proprietà per assurdo, facendo ricorso alla classica tecnica del \textit{cut-and-paste}.

Sia $R^*$ una soluzione ottima per il sottoproblema generico $(i, b)$, ovvero una selezione ammissibile di giocatori scelti tra i primi $i$ ruoli che rispetta il vincolo di bilancio residuo $\le b$ e consegue il massimo valore cumulativo $m[i, b]$. Rispetto alla decisione assunta per l'$i$-esimo ruolo, si presentano due casi mutualmente esclusivi ed esaustivi:

\begin{enumerate}
    \item \textbf{Caso 1: Nessun giocatore svincolato viene ingaggiato per il ruolo $i$.} \\
    In questo caso la rosa $R^*$ è interamente composta da giocatori appartenenti ai primi $i-1$ ruoli con spesa cumulativa non superiore a $b$. Di conseguenza, $R^*$ costituisce una soluzione ammissibile per il sottoproblema $(i-1, b)$. \\
    \textit{Dimostrazione}: se per assurdo $R^*$ non fosse ottima per $(i-1, b)$, allora esisterebbe una diversa configurazione ammissibile $R'$ per i primi $i-1$ ruoli, con costo $\le b$, tale che $\text{WAR}(R') > \text{WAR}(R^*)$. Poiché $R'$ non alloca alcun giocatore al ruolo $i$, essa rappresenterebbe anche una soluzione ammissibile per $(i, b)$ con rendimento $\text{WAR}(R') > \text{WAR}(R^*)$, contraddicendo l'ipotesi di ottimalità globale di $R^*$ su $(i, b)$.

    \item \textbf{Caso 2: Viene ingaggiato il candidato $p = \text{players}[i-1, j^*]$ per il ruolo $i$, avente costo scalato $c_{i, j^*} \le b$ e valore $p.war$.} \\
    Rimuovendo $p$ dalla soluzione ottima, consideriamo la sottorosa $R' = R^* \setminus \{p\}$. Poiché $R^*$ rispetta il vincolo di bilancio complessivo $\le b$, la sottorosa $R'$ impegna una quota di spesa non superiore a $b - c_{i, j^*}$ sui primi $i-1$ ruoli, qualificandosi come soluzione ammissibile per il sottoproblema $(i-1, b - c_{i, j^*})$, con rendimento pari a $\text{WAR}(R') = \text{WAR}(R^*) - p.war$. \\
    \textit{Dimostrazione}: supponiamo per assurdo che $R'$ non sia una soluzione ottima per il sottoproblema $(i-1, b - c_{i, j^*})$. Allora deve esistere una rosa ammissibile alternativa $R''$ per i primi $i-1$ ruoli, con spesa $\le b - c_{i, j^*}$, tale che:
    \[
    \text{WAR}(R'') > \text{WAR}(R')
    \]
    Possiamo quindi costruire una nuova rosa $R_{\text{new}} = R'' \cup \{p\}$ per i primi $i$ ruoli. Verifichiamo che $R_{\text{new}}$ sia ammissibile per il sottoproblema $(i, b)$:
    \begin{itemize}
        \item \textit{Vincolo di rosa}: $R''$ include al più un giocatore per ciascuno dei primi $i-1$ ruoli; aggiungendo $p$ (che ricopre il ruolo $i$), ogni posizione tra $1$ e $i$ ospita al massimo un elemento.
        \item \textit{Vincolo di bilancio}: il costo complessivo di $R''$ è $\le b - c_{i, j^*}$. Aggiungendo il costo del giocatore $p$, la spesa totale diventa $\le (b - c_{i, j^*}) + c_{i, j^*} = b \le B$.
    \end{itemize}
    Calcolando il rendimento complessivo della nuova configurazione:
    \[
    \text{WAR}(R_{\text{new}}) = \text{WAR}(R'') + p.war > \text{WAR}(R') + p.war = \text{WAR}(R^*)
    \]
    Abbiamo così ottenuto una rosa ammissibile per $(i, b)$ avente un rendimento strettamente superiore a quello di $R^*$, il che contraddice l'assunto di ottimalità di $R^*$. Ne consegue che la sottorosa $R'$ deve necessariamente essere una soluzione ottima per il sottoproblema $(i-1, b - c_{i, j^*})$.
\end{enumerate}
In entrambi i casi, la soluzione ottima globale incorpora soluzioni ottime per i sottoproblemi correlati. La proprietà di sottostruttura ottima è pertanto verificata.

\subsubsection{Sovrapposizione dei Sottoproblemi}
La seconda proprietà richiesta per l'efficacia del paradigma dinamico riguarda la presenza di sottoproblemi comuni ricorsivamente condivisi:
\begin{itemize}
    \item \textbf{Spazio totale dei sottoproblemi distinti:} ogni sottoproblema è univocamente individuato dalla coppia di indici $(i, b)$, con $0 \le i \le N$ e $0 \le b \le B$. Il numero totale di sottoproblemi distinti è strettamente limitato e pari a:
    \[
    (N + 1) \times (B + 1)
    \]
    \item \textbf{Inefficienza di un approccio ricorsivo ingenuo:} se si tentasse di risolvere il problema mediante un algoritmo puramente ricorsivo (o basato su divide-et-impera standard senza memorizzazione), l'algoritmo tenterebbe di calcolare tutte le combinazioni possibili a ogni invocazione. Di conseguenza, a ogni livello $i$ si aprirebbero fino a $P + 1$ rami di scelta (non ingaggiare nessuno o selezionare uno tra i $P$ candidati). L'albero di ricorsione genererebbe complessivamente $\Omega((P + 1)^N)$ chiamate, crescendo in modo esponenziale rispetto al numero di ruoli da coprire.
    \item \textbf{Sottoproblemi comuni calcolati ripetutamente:} in questo problema scelte di ingaggio diverse in ruoli differenti possono portare al medesimo ruolo $i$ con lo stesso budget residuo $b$. Ad esempio, ingaggiare un giocatore da 2 milioni al ruolo 1 e nessuno al ruolo 2 produce esattamente lo stesso budget residuo che ingaggiare due giocatori da 1 milione ai ruoli 1 e 2. Un algoritmo privo di memoria risolverebbe ciascuno stato comune $(i, b)$ un numero esponenziale di volte.
\end{itemize}
La programmazione dinamica risolve ciascuno dei sottoproblemi una sola volta, memorizzandone l'esito nella tabella $m[i, b]$ e garantendo che la consultazione avvenga in tempo costante $\Theta(1)$ nei passi successivi.

\subsection{Equazione di Ricorrenza}
In virtù delle proprietà dimostrate, il valore della soluzione ottima $m[i, b]$ è calcolabile mediante la seguente equazione di ricorrenza:
\[
\begin{cases} 
0 & \text{se } i = 0 \text{ oppure } b = 0 \\
\max \left( m[i - 1, b], \; \max\limits_{\substack{0 \le j < P \\ c_{i, j} \le b}} \{ m[i - 1, b - c_{i, j}] + \text{players}[i-1, j].war \} \right) & \text{se } i > 0 \text{ e } b > 0
\end{cases}
\]
Nel passo ricorsivo:
\begin{itemize}
    \item il primo termine $m[i - 1, b]$ rappresenta la scelta di non ingaggiare alcun giocatore per il ruolo $i$ (mantenendo inalterato il WAR conseguito sui primi $i-1$ ruoli senza consumare budget aggiuntivo);
    \item il secondo termine valuta l'opzione di ingaggiare il miglior candidato $j$ tra i $P$ disponibili per il ruolo $i$, a patto che il suo costo scalato $c_{i, j}$ non ecceda il budget residuo $b$.
\end{itemize}

\subsection{Pseudocodici degli Algoritmi Fondamentali}
In conformità a quanto richiesto dalle specifiche progettuali, presentiamo di seguito gli pseudocodici formali per i due algoritmi risolutivi fondamentali:
\begin{enumerate}
    \item \textbf{Calcolo del valore ottimo}: un algoritmo iterativo \textit{bottom-up} che calcola il massimo WAR realizzabile memorizzando i risultati parziali nella tabella $m$ e tracciando le scelte effettuate nella tabella ausiliaria $s$.
    \item \textbf{Ricostruzione della soluzione ottima}: un algoritmo che sfrutta le informazioni memorizzate nella tabella $s$ per estrarre la composizione esatta della rosa ottima e il costo cumulativo finale.
\end{enumerate}

\subsubsection{Algoritmo 1: Calcolo del Valore Ottimo e Tabella delle Scelte}
L'Algoritmo~\ref{alg:compute_optimal_war} opera secondo una strategia di popolamento tabellare progressivo. Le righe della tabella corrispondono ai ruoli considerati ($i = 0, \dots, N$), mentre le colonne corrispondono ai valori discreti di budget ammissibili ($b = 0, \dots, B$).

\begin{algorithm}[H]
\caption{Calcolo del WAR Ottimale e della Tabella delle Scelte}
\label{alg:compute_optimal_war}
\begin{algorithmic}[1]
\Procedure{Compute-Optimal-War}{$N, X, P, \mathit{players}$}
    \State $\text{COST\_UNIT} \gets 100\,000$
    \State $B \gets \lfloor X / \text{COST\_UNIT} \rfloor$
    \State alloca la matrice $m[0 \dots N, 0 \dots B]$ inizializzata a $0.0$
    \State alloca la matrice $s[0 \dots N, 0 \dots B]$ inizializzata a $-1$
    \For{$i \gets 1$ \textbf{to} $N$} \Comment{Considera ciascun ruolo difensivo}
        \State $m[i, 0 \dots B] \gets m[i - 1, 0 \dots B]$ \Comment{Scelta base: mantieni il rendimento precedente}
        \For{$j \gets 0$ \textbf{to} $P - 1$} \Comment{Esamina ciascun candidato disponibile per il ruolo $i$}
            \State $p \gets \mathit{players}[i - 1, j]$
            \State $c \gets \lfloor p.\mathit{cost} / \text{COST\_UNIT} \rfloor$
            \For{$b \gets c$ \textbf{to} $B$} \Comment{Aggiorna gli stati in cui il budget residuo è sufficiente}
                \State $\mathit{war} \gets m[i - 1, b - c] + p.\mathit{war}$
                \If{$\mathit{war} > m[i, b]$}
                    \State $m[i, b] \gets \mathit{war}$
                    \State $s[i, b] \gets j$ \Comment{Memorizza l'indice del giocatore selezionato}
                \EndIf
            \EndFor
        \EndFor
    \EndFor
    \State $\mathit{optimal\_war} \gets m[N, B]$
    \State \textbf{return} $\mathit{optimal\_war}, s$
\EndProcedure
\end{algorithmic}
\end{algorithm}

\subsubsection{Algoritmo 2: Ricostruzione della Soluzione Ottima}
L'Algoritmo~\ref{alg:reconstruct_optimal_roster} ricostruisce l'elenco dei giocatori ingaggiati effettuando un percorso a ritroso (\textit{backtracking}) a partire dalla cella di arrivo $(N, B)$ della matrice delle decisioni $s$.

\begin{algorithm}[H]
\caption{Ricostruzione della Rosa Ottima}
\label{alg:reconstruct_optimal_roster}
\begin{algorithmic}[1]
\Procedure{Reconstruct-Optimal-Roster}{$N, X, \mathit{players}, s$}
    \State $\text{COST\_UNIT} \gets 100\,000$
    \State $b \gets \lfloor X / \text{COST\_UNIT} \rfloor$
    \State $\mathit{signedPlayers} \gets$ nuova lista vuota
    \State $\mathit{totalCost} \gets 0$
    \For{$i \gets N$ \textbf{downto} $1$} \Comment{Ripercorre le decisioni dall'ultimo ruolo al primo}
        \State $\mathit{index} \gets s[i, b]$
        \If{$\mathit{index} \neq -1$} \Comment{Un giocatore svincolato è stato effettivamente ingaggiato}
            \State $p \gets \mathit{players}[i - 1, \mathit{index}]$
            \State inserisci $p$ in $\mathit{signedPlayers}$
            \State $\mathit{totalCost} \gets \mathit{totalCost} + p.\mathit{cost}$
            \State $b \gets b - \lfloor p.\mathit{cost} / \text{COST\_UNIT} \rfloor$ \Comment{Scala il costo dal budget residuo}
        \EndIf
    \EndFor
    \State \textbf{return} $\mathit{signedPlayers}, \mathit{totalCost}$
\EndProcedure
\end{algorithmic}
\end{algorithm}

\subsection{Analisi Asintotica Teorica del Primo Algoritmo}
L'analisi asintotica formale dell'Algoritmo~\ref{alg:compute_optimal_war} viene condotta valutando separatamente la complessità temporale, l'occupazione di memoria (complessità spaziale) e la classe di complessità teorica dell'approccio proposto.

\subsubsection{Complessità Temporale}
Il tempo di esecuzione dell'algoritmo è determinato dalle seguenti fasi:
\begin{enumerate}
    \item \textbf{Fase di inizializzazione}: l'allocazione e l'azzeramento delle due matrici $m$ ed $s$ di dimensioni $(N + 1) \times (B + 1)$ richiedono:
    \[
    T_{\text{init}} = \Theta(N \cdot B)
    \]
    \item \textbf{Fase di calcolo iterativo (cicli principali)}:
    \begin{itemize}
        \item Il ciclo più esterno (linea 6) viene eseguito esattamente $N$ volte, per $i = 1, \dots, N$.
        \item All'inizio di ogni iterazione del ruolo $i$, la copia del vettore base dallo stage precedente $m[i-1, 0 \dots B]$ (linea 7) comporta $B + 1$ assegnamenti, pari a $\Theta(B)$ operazioni elementari.
        \item Il secondo ciclo (linea 8) itera sui $P$ candidati disponibili per il ruolo $i$.
        \item Per ogni candidato $j$, il ciclo più interno sul budget (linea 11) viene eseguito per $b$ che varia da $c$ a $B$, compiendo esattamente $B - c + 1$ iterazioni. Poiché nel caso peggiore $c \ge 0$, tale ciclo compie al più $B + 1$ iterazioni.
        \item All'interno del ciclo più interno, il calcolo del valore candidato, il confronto con il valore corrente e l'eventuale aggiornamento delle matrici $m[i, b]$ e $s[i, b]$ (linee 12--15) richiedono tempo costante $\Theta(1)$.
    \end{itemize}
\end{enumerate}

Sommando i contributi per tutti i ruoli e candidati, il numero complessivo di operazioni elementari $T(N, B, P)$ è formalmente delimitato da:
\[
T(N, B, P) = \Theta(N \cdot B) + \sum_{i=1}^N \left( \Theta(B) + \sum_{j=0}^{P-1} \sum_{b=c_{i,j}}^B \Theta(1) \right)
\]
Poiché $0 \le c_{i,j} \le B$, il termine $\sum_{b=c_{i,j}}^B \Theta(1) \le (B + 1) \cdot \Theta(1) = \Theta(B)$. Dunque:
\[
T(N, B, P) = \Theta(N \cdot B) + \sum_{i=1}^N \Big( \Theta(B) + P \cdot \Theta(B) \Big) = \Theta(N \cdot B) + \Theta(N \cdot B \cdot P) = \Theta(N \cdot B \cdot P)
\]
Nel caso generale in cui il numero di candidati $P_i$ vari ruolo per ruolo, la complessità temporale si esprime più precisamente come:
\[
T\left(N, B, \{P_i\}\right) = \Theta\left(B \sum_{i=1}^N P_i\right)
\]
che, assumendo $P_i = P$ costante per ogni posizione, coincide esattamente con $\Theta(N \cdot B \cdot P)$.

\subsubsection{Occupazione di Memoria (Complessità Spaziale)}
L'occupazione di memoria dell'algoritmo è dominata dalle strutture dati allocate per il mantenimento dello stato dinamico:
\begin{itemize}
    \item La matrice dei valori $m$ ha dimensioni $(N + 1) \times (B + 1)$. Ciascuna cella memorizza un numero in virgola mobile a precisione doppia (\texttt{float64}, 8 byte) che rappresenta il WAR cumulativo ottimo parziale.
    \item La matrice delle decisioni $s$ ha anch'essa dimensioni $(N + 1) \times (B + 1)$. Ciascuna cella memorizza un indice intero (\texttt{int}, 4 o 8 byte) corrispondente al candidato prescelto (oppure $-1$ per indicare l'assenza di ingaggi).
    \item Le variabili scalari ausiliarie ($N, X, B, P, i, j, b, c, \mathit{war}$) richiedono una quantità di memoria indipendente dalla taglia dell'istanza, pari a $O(1)$.
\end{itemize}
La complessità spaziale complessiva per l'esecuzione del primo algoritmo è pertanto:
\[
S(N, B) = 2 \times (N + 1) \times (B + 1) \cdot O(1) = \Theta(N \cdot B)
\]
Si noti che l'occupazione di memoria dell'Algoritmo~\ref{alg:compute_optimal_war} è completamente indipendente dal numero di candidati per ruolo $P$, dipendendo unicamente dal numero di ruoli $N$ e dal budget discretizzato $B$.

\subsubsection{Discussione sulla Pseudo-Polinomialità}
Una considerazione teorica fondamentale, richiesta dall'analisi della complessità computazionale, riguarda la natura della complessità riscontrata:
\begin{quote}
\textit{L'algoritmo proposto presenta una complessità temporale $\Theta(N \cdot B \cdot P)$ e spaziale $\Theta(N \cdot B)$ che si qualificano come \textbf{pseudo-polinomiali}.}
\end{quote}

Formalmente, la teoria della complessità misura l'efficienza di un algoritmo non rispetto al valore numerico delle grandezze in ingresso, bensì rispetto alla \textbf{lunghezza della codifica in bit dell'input} (detta dimensione dell'istanza). 

Nel nostro problema, il parametro $B$ rappresenta un valore numerico la cui rappresentazione binaria richiede un numero di bit pari a:
\[
L_B = \lfloor \log_2 B \rfloor + 1
\]
Di conseguenza, il valore numerico $B$ è esponenziale rispetto al numero di bit necessari per memorizzarlo:
\[
B = \Theta(2^{L_B})
\]
Esprimendo la complessità temporale in funzione della dimensione dell'input in bit, si ottiene:
\[
T(N, L_B, P) = \Theta(N \cdot 2^{L_B} \cdot P)
\]
L'algoritmo manifesta una dipendenza esponenziale dalla lunghezza della codifica del budget $B$. Pertanto, non si tratta di un algoritmo polinomiale in senso stretto. 

Questa proprietà colloca il Multiple-Choice Knapsack Problem (MCKP), analogamente al classico problema dello zaino 0/1 (Knapsack) e al Subset Sum, nella classe dei problemi \textbf{NP-hard in senso debole} (\textit{weakly NP-hard}):
\begin{itemize}
    \item l'algoritmo risulta estremamente efficiente e trattabile nella pratica fintanto che il budget scalato $B$ si mantiene di entità moderata;
    \item la complessità computazionale e l'occupazione di memoria degradano rapidamente verso l'intrattabilità qualora il budget assuma ordini di grandezza arbitrariamente elevati o qualora si scelga una scala monetaria infinitesima ($\text{COST\_UNIT} \to 1$).
\end{itemize}
